\documentclass[10pt,a4paper]{amsart}
\usepackage{amsmath,amssymb,amsthm,graphicx}
\usepackage{verbatim}
\newenvironment{reference}{\comment}{\endcomment}
%\newenvironment{reference}{}{}
\newenvironment{slogan}{\comment}{\endcomment}
\newenvironment{history}{\comment}{\endcomment}

% For commutative diagrams we use Xy-pic
\usepackage[all]{xy}

% We use 2cell for 2-commutative diagrams.
\xyoption{2cell}
% Theorem environments.
%
\theoremstyle{plain}
\newtheorem{theorem}[subsection]{Theorem}
\newtheorem{proposition}[subsection]{Proposition}
\newtheorem{lemma}[subsection]{Lemma}

\theoremstyle{definition}
\newtheorem{definition}[subsection]{Definition}
\newtheorem{example}[subsection]{Example}
\newtheorem{exercise}[subsection]{Exercise}
\newtheorem{situation}[subsection]{Situation}

\theoremstyle{remark}
\newtheorem{remark}[subsection]{Remark}
\newtheorem{remarks}[subsection]{Remarks}

\numberwithin{equation}{subsection}

% Macros
%
\def\lim{\mathop{\mathrm{lim}}\nolimits}
\def\colim{\mathop{\mathrm{colim}}\nolimits}
\def\Spec{\mathop{\mathrm{Spec}}}
\def\Hom{\mathop{\mathrm{Hom}}\nolimits}
\def\Ext{\mathop{\mathrm{Ext}}\nolimits}
\def\SheafHom{\mathop{\mathcal{H}\!\mathit{om}}\nolimits}
\def\SheafExt{\mathop{\mathcal{E}\!\mathit{xt}}\nolimits}
\def\Sch{\mathit{Sch}}
\def\Mor{\mathop{\mathrm{Mor}}\nolimits}
\def\Ob{\mathop{\mathrm{Ob}}\nolimits}
\def\Sh{\mathop{\mathit{Sh}}\nolimits}
\def\NL{\mathop{N\!L}\nolimits}
\def\CH{\mathop{\mathrm{CH}}\nolimits}
\def\proetale{{pro\text{-}\acute{e}tale}}
\def\etale{{\acute{e}tale}}
\def\QCoh{\mathit{QCoh}}
\def\Ker{\mathop{\mathrm{Ker}}}
\def\Im{\mathop{\mathrm{Im}}}
\def\Coker{\mathop{\mathrm{Coker}}}
\def\Coim{\mathop{\mathrm{Coim}}}

% Boxtimes
%
\DeclareMathSymbol{\boxtimes}{\mathbin}{AMSa}{"02}

%
% Macros for moduli stacks/spaces
%
\def\QCohstack{\mathcal{QC}\!\mathit{oh}}
\def\Cohstack{\mathcal{C}\!\mathit{oh}}
\def\Spacesstack{\mathcal{S}\!\mathit{paces}}
\def\Quotfunctor{\mathrm{Quot}}
\def\Hilbfunctor{\mathrm{Hilb}}
\def\Curvesstack{\mathcal{C}\!\mathit{urves}}
\def\Polarizedstack{\mathcal{P}\!\mathit{olarized}}
\def\Complexesstack{\mathcal{C}\!\mathit{omplexes}}
% \Pic is the operator that assigns to X its picard group, usage \Pic(X)
% \Picardstack_{X/B} denotes the Picard stack of X over B
% \Picardfunctor_{X/B} denotes the Picard functor of X over B
\def\Pic{\mathop{\mathrm{Pic}}\nolimits}
\def\Picardstack{\mathcal{P}\!\mathit{ic}}
\def\Picardfunctor{\mathrm{Pic}}
\def\Deformationcategory{\mathcal{D}\!\mathit{ef}}
\usepackage{cleveref}
\crefname{originalreference}{original reference}{original references}
\Crefname{originalreference}{Original reference}{Original references}
\makeatletter
\newlabel{lemma-etale-image-open}{{06NF}{source}{}{sourcepointer-1}{}}
\newlabel{lemma-etale-image-open@cref}{{[originalreference][0][]06NF}{[1][0][]source}{}{}{}}
\makeatother
\crefname{originalreference}{original reference}{original references}
\Crefname{originalreference}{Original reference}{Original references}
\makeatletter
\newlabel{lemma-union-of-quasi-compact}{{03FY}{source}{}{sourcepointer-2}{}}
\newlabel{lemma-union-of-quasi-compact@cref}{{[originalreference][0][]03FY}{[1][0][]source}{}{}{}}
\makeatother
\crefname{originalsection}{original section}{original sections}
\Crefname{originalsection}{Original section}{Original sections}
\makeatletter
\newlabel{schemes-section-fibre-products}{{01JO}{source}{}{sourcepointer-3}{}}
\newlabel{schemes-section-fibre-products@cref}{{[originalsection][0][]01JO}{[1][0][]source}{}{}{}}
\makeatother
\crefname{originallemma}{original lemma}{original lemmas}
\Crefname{originallemma}{Original lemma}{Original lemmas}
\makeatletter
\newlabel{spaces-lemma-composition-representable-transformations-property}{{02WK}{source}{}{sourcepointer-4}{}}
\newlabel{spaces-lemma-composition-representable-transformations-property@cref}{{[originallemma][0][]02WK}{[1][0][]source}{}{}{}}
\makeatother
\crefname{originallemma}{original lemma}{original lemmas}
\Crefname{originallemma}{Original lemma}{Original lemmas}
\makeatletter
\newlabel{spaces-lemma-product-representable-transformations-property}{{02WM}{source}{}{sourcepointer-5}{}}
\newlabel{spaces-lemma-product-representable-transformations-property@cref}{{[originallemma][0][]02WM}{[1][0][]source}{}{}{}}
\makeatother
\crefname{originallemma}{original lemma}{original lemmas}
\Crefname{originallemma}{Original lemma}{Original lemmas}
\makeatletter
\newlabel{spaces-lemma-representable-morphisms-spaces-property}{{02WZ}{source}{}{sourcepointer-6}{}}
\newlabel{spaces-lemma-representable-morphisms-spaces-property@cref}{{[originallemma][0][]02WZ}{[1][0][]source}{}{}{}}
\makeatother
\title{Characterizations of being Zariski locally quasi-separated through covers by affine schemes}
\author{The Stacks Project Authors}
\date{stacks-project:spaces-properties.tex:local-snapshot-20261012; GNU FDL 1.2 or later}
\begin{document}
\maketitle
\noindent\textit{Editorial note.} All delivered mathematical text is the author's own, copied from the pinned Stacks Project chapter file spaces-properties.tex; no statement or proof step was rewritten, reordered or removed.\par\medskip
\noindent\textit{Editorial note.} The counted proof is the author's complete argument: the refinement of a cover into one with affine pieces over affine opens of the base, the deduction of quasi-compactness of the diagonal from the criterion on a surjective etale cover, and the reverse implication using the existence lemma and the quasi-compactness of the fibre products over the pieces.\par\medskip
\noindent\textit{Editorial note.} The label lemma-quasi-separated-quasi-compact-pieces names the counted statement, and the proof environment immediately following it in the source is delivered verbatim as its proof.\par\medskip
\noindent\textit{Editorial note.} The counted statement is carried by the author's own lemma environment in the chapter source.\par\medskip
\noindent\textit{Editorial note.} Pre-open gate: python operations/agent-authoring/tools/precard\_gate.py --source spaces-properties.tex --label lemma-quasi-separated-quasi-compact-pieces --start 1064 --end 1149 --local-id 28043 returned RESULT=OK (L28043 free, occupancy FREE, public interval clear). Candidate taken from sources/stacks-spaces-properties/candidates-halfA.json.\par\medskip
\noindent\textit{Editorial note.} The statement and proof are copied from the pinned chapter file; the cited inputs are kept as external pointers with their canonical Stacks tags.\par\medskip
\noindent\textit{Allowed standard prerequisites (editorial).}\begin{itemize}
\item Openness of the image of an etale morphism of algebraic spaces. Same chapter, the lemma that the image of an etale morphism of algebraic spaces is open. Supplies the open subspace X\_ij corresponding to the image of the map from a piece of the affine cover.
\item Quasi-compactness of a representable morphism tested on a surjective etale cover. Spaces, the lemma reducing a property of a representable morphism to a surjective etale cover by a scheme. Supplies the criteria by which quasi-compactness of the diagonal and of the morphisms U\_i -> X\_i is concluded.
\item Existence of a Zariski cover by pieces mapping into affine opens with affine etale covers. Same chapter, the lemma producing a Zariski open covering by pieces mapping into affine opens of the base and admitting surjective etale maps from affine schemes. Supplies the covering used in the reduction from the Zariski local case.
\end{itemize}\medskip
\ifdefined\hypertarget\hypertarget{sourcepointer-1}{}\fi
\noindent\textit{Source reference.} Original reference 06NF: sources/stacks-spaces-properties/pinned-original-source.tex, source lines 753--753. Prerequisite pointer invoked at proof line(s) 1093; it names a standard fact of the same chapter that the counted argument uses, and every step of the counted proof built on it is retained verbatim in the delivered body.\par\medskip
\ifdefined\hypertarget\hypertarget{sourcepointer-2}{}\fi
\noindent\textit{Source reference.} Original reference 03FY: sources/stacks-spaces-properties/pinned-original-source.tex, source lines 943--943. Prerequisite pointer invoked at proof line(s) 1125; it names a standard fact of the same chapter that the counted argument uses, and every step of the counted proof built on it is retained verbatim in the delivered body.\par\medskip
\ifdefined\hypertarget\hypertarget{sourcepointer-3}{}\fi
\noindent\textit{Source reference.} Original section 01JO: sources/stacks-chow/tags.txt, source lines 2021--2021. Prerequisite pointer invoked at proof line(s) 1115; it names a standard fact of a sibling chapter that the counted argument uses, and every step of the counted proof built on it is retained verbatim in the delivered body.\par\medskip
\ifdefined\hypertarget\hypertarget{sourcepointer-4}{}\fi
\noindent\textit{Source reference.} Original lemma 02WK: sources/stacks-chow/tags.txt, source lines 3807--3807. Prerequisite pointer invoked at proof line(s) 1090; it names a standard fact of a sibling chapter that the counted argument uses, and every step of the counted proof built on it is retained verbatim in the delivered body.\par\medskip
\ifdefined\hypertarget\hypertarget{sourcepointer-5}{}\fi
\noindent\textit{Source reference.} Original lemma 02WM: sources/stacks-chow/tags.txt, source lines 3809--3809. Prerequisite pointer invoked at proof line(s) 1113; it names a standard fact of a sibling chapter that the counted argument uses, and every step of the counted proof built on it is retained verbatim in the delivered body.\par\medskip
\ifdefined\hypertarget\hypertarget{sourcepointer-6}{}\fi
\noindent\textit{Source reference.} Original lemma 02WZ: sources/stacks-chow/tags.txt, source lines 3822--3822. Prerequisite pointer invoked at proof line(s) 1119; it names a standard fact of a sibling chapter that the counted argument uses, and every step of the counted proof built on it is retained verbatim in the delivered body.\par\medskip
\begin{lemma}
\label{lemma-quasi-separated-quasi-compact-pieces}
Let $S$ be a scheme.
Let $X$ be an algebraic space over $S$.
The following are equivalent
\begin{enumerate}
\item $X$ is Zariski locally quasi-separated over $S$,
\item $X$ is Zariski locally quasi-separated,
\item there exists a Zariski open covering $X = \bigcup X_i$
such that for each $i$ there exists an affine scheme
$U_i$ and a quasi-compact surjective \'etale
morphism $U_i \to X_i$, and
\item there exists a Zariski open covering $X = \bigcup X_i$
such that for each $i$ there exists an affine scheme
$U_i$ which maps into an affine open of $S$ and a quasi-compact
surjective \'etale morphism $U_i \to X_i$.
\end{enumerate}
\end{lemma}

\begin{proof}
Assume  $U_i \to X_i \subset X$ are as in (3). To prove (4)
choose for each $i$ a finite affine open covering $U_i =
U_{i1} \cup \ldots \cup U_{in_i}$ such that each $U_{ij}$ maps
into an affine open of $S$. The compositions
$U_{ij} \to U_i \to X_i$ are \'etale and quasi-compact (see
Spaces, Lemma
\ref{spaces-lemma-composition-representable-transformations-property}).
Let $X_{ij} \subset X_i$ be the open subspace corresponding to
the image of $|U_{ij}| \to |X_i|$, see
Lemma \ref{lemma-etale-image-open}.
Note that $U_{ij} \to X_{ij}$ is quasi-compact as $X_{ij} \subset X_i$
is a monomorphism and as $U_{ij} \to X$ is quasi-compact.
Then $X = \bigcup X_{ij}$ is a covering as in (4).
The implication (4) $\Rightarrow$ (3) is immediate.

\medskip\noindent
Assume (4). To show that $X$ is Zariski locally quasi-separated over $S$
it suffices to show that $X_i$ is quasi-separated over $S$.
Hence we may assume there exists an affine scheme $U$ mapping into
an affine open of $S$ and a quasi-compact surjective \'etale
morphism $U \to X$. Consider the fibre product square
$$
\xymatrix{
U \times_X U \ar[r] \ar[d] & U \times_S U \ar[d] \\
X \ar[r]^-{\Delta_{X/S}} & X \times_S X
}
$$
The right vertical arrow is surjective \'etale (see
Spaces, Lemma
\ref{spaces-lemma-product-representable-transformations-property})
and $U \times_S U$ is affine (as $U$ maps into an affine open of $S$, see
Schemes, Section \ref{schemes-section-fibre-products}),
and $U \times_X U$ is quasi-compact
because the projection $U \times_X U \to U$ is quasi-compact as a
base change of $U \to X$. It follows from
Spaces, Lemma \ref{spaces-lemma-representable-morphisms-spaces-property}
that $\Delta_{X/S}$ is quasi-compact as desired.

\medskip\noindent
Assume (1). To prove (3) there is an immediate reduction to the case
where $X$ is quasi-separated over $S$. By
Lemma \ref{lemma-union-of-quasi-compact}
we can find a Zariski open covering $X = \bigcup X_i$ such that each
$X_i$ maps into an affine open of $S$, and such that there exist affine
schemes $U_i$ and surjective \'etale morphisms $U_i \to X_i$.
Since $U_i \to S$ maps into an affine open of $S$ we see that
$U_i \times_S U_i$ is affine, see
Schemes, Section \ref{schemes-section-fibre-products}.
As $X$ is quasi-separated over $S$, the morphisms
$$
R_i = U_i \times_{X_i} U_i = U_i \times_X U_i
\longrightarrow
U_i \times_S U_i
$$
as base changes of $\Delta_{X/S}$ are quasi-compact. Hence we conclude
that $R_i$ is a quasi-compact scheme. This in turn implies that each
projection $R_i \to U_i$ is quasi-compact. Hence, applying
Spaces, Lemma \ref{spaces-lemma-representable-morphisms-spaces-property}
to the covering $U_i \to X_i$ and the morphism $U_i \to X_i$
we conclude that the morphisms $U_i \to X_i$ are quasi-compact as desired.

\medskip\noindent
At this point we see that (1), (3), and (4) are equivalent. Since (3) does
not refer to the base scheme we conclude that these are also equivalent
with (2).
\end{proof}

\end{document}
